% RESEARCH_DRAFT_SHA256: 2b1219457d1a00a914934eb5706642ac4c40aca33525023952d01b3073afdd2a
% Layout and source-reference wrappers expanded for the web renderer.
\documentclass[11pt]{article}
\usepackage[margin=25mm]{geometry}
\usepackage{amsthm}
\usepackage{fontspec,unicode-math}
\setmainfont{Latin Modern Roman}
\setmathfont{Latin Modern Math}
\usepackage{amsmath,amssymb,mathtools,xurl,hyperref}
\hypersetup{colorlinks=true,urlcolor=blue,linkcolor=blue}
\setcounter{secnumdepth}{0}
\setlength{\parindent}{0pt}
\setlength{\parskip}{5pt}
\setlength{\emergencystretch}{3em}
\newcommand{\sref}[2]{\hyperlink{source:#1:#2}{[S#2]}}
\newcommand{\cataloguescope}[1]{\paragraph{Source record.}#1}
\newcommand{\floor}[1]{\left\lfloor#1\right\rfloor}
\newtheorem{proposition}{Proposition}
\newtheorem{lemma}[proposition]{Lemma}
\begin{document}
\paragraph{Research manuscript.} OpenAI \texttt{gpt-6-astra}, reasoning \texttt{ultra}; 11 October 2026.
\paragraph{Preserved source body.} The complete catalogue mathematical text and sources follow with layout and source links expanded, including source qualifications and its historical preparation notes. New research begins at the section headed Research attempt.

% UnsolvedMath ID 30002626; source catalogue code OWR-13105-003
\section{Optimal Fair Multiparty Coin Flipping}\label{30002626}
{\small UnsolvedMath ID 30002626 · Computer Science}
\subsection{Definitions and mathematical statement}

\paragraph{Shared notation.}Standard number systems are $\mathbb N,\mathbb Z,\mathbb Q,\mathbb R,\mathbb C$. Specialized notation and index ranges follow the cited source and the definitions below.


In multiparty coin flipping, participants aim to produce a common random bit despite adversarial participants. Fairness restricts the adversary\textquotesingle s ability to bias output or obtain an outcome while preventing honest participants from doing so.
\paragraph{Record scope and qualifications.}
The catalogue title is retained as metadata. Named certificates, contracts, bounds or frameworks in that title are not asserted as results of the original source extract.
The meaning of $m$, overhead measure, communication model, and precise fairness notion are absent from the supplied question.

Source-record limitation: The variable m, overhead measure, corruption model, and precise fairness requirement are not specified.
\paragraph{Statement recorded in the supplied source.}
Does there exist a fair three-party coin-flipping protocol without the $O(\log^2 m)$ factor? More generally, does there exist a fair protocol for any number of parties against any number of corrupted parties? \href{https://doi.org/10.4171/owr/2014/35}{[S2]}
\paragraph{Source status.}
The supplied catalogue labels this source record open; that metadata was checked against the supplied snapshot on 2026-10-11. The label is preserved for provenance and is not an independent verification of present open status or mathematical correctness. The catalogue literature assessment is dated 2026-08-21.
\subsection{Short English statement}
Remove the stated overhead in fair three-party coin flipping and determine what fairness is achievable with more parties and corruptions.
\subsection{Sources}
\begin{itemize}
\item[S1] Ulam AI and the UnsolvedMath Contributors, UnsolvedMath: A Curated Collection of Open Mathematics Problems, record 30002626 (OWR-13105-003). Catalogue curation and metadata: CC BY 4.0. The source extract is retained; definitions, English summaries and qualifications are editorial additions.
\url{https://huggingface.co/datasets/ulamai/UnsolvedMath}.
\item[S2] Cryptography (2015), Oberwolfach Reports, DOI 10.4171/owr/2014/35, pp. 1933--1990
\url{https://doi.org/10.4171/owr/2014/35}.
\item[S3] Haitner--Tsfadia, fair coin flipping
\url{https://epubs.siam.org/doi/10.1137/15M1009147}.
\item[S4] 2025 many-party coin-flipping bounds
\url{https://link.springer.com/article/10.1007/s00145-025-09561-6}.
\item Source-specific Oberwolfach Reports licences, including CC BY-SA 4.0 where applicable, are retained; see THIRD\_PARTY\_NOTICES.md.
\end{itemize}
\paragraph{Preparation and review.}
The original source and UnsolvedMath contributors retain attribution for the source extract. Codex (GPT-6; exact serving revision not exposed) assisted with definitions, English summaries, scope notes and formatting. No new human mathematical authorship or independently verified solution is asserted. Permanent source ID 30002626 and code OWR-13105-003 checked against the complete supplied registry and current repository. Distinct upstream records, including repeated questions and broad or damaged extracts, are retained for complete coverage of the source CS category. A repeated record is not a new independent conjecture. Complete mathematical-equivalence and current-literature review has not been performed. Preparation and final review follow contribution prompts C and D, with source gaps disclosed. Structural and rendering checks do not establish a new theorem or open status.
\label{end:30002626}

\section{Research attempt}
\subsection{Literature, source context and scope}
The full original catalogue body and source qualifications are preserved above. The relevant primary workshop passages were read in the cached publisher report. Historical claims and subsequent attributed results are distinguished from the mathematical work below. A restricted example, standard lemma or known theorem is not counted as new progress on the entire original target. Missing assumptions and repeated or rhetorical extracts are not silently turned into different solved problems.

\subsection{A visible preliminary bit permits selective-abort bias}
The primary source specifies the communication model and uses $m$ for the number of rounds; its three-party construction and target bias are preserved above. We test a much simpler attempted building block, not that construction or its lower bounds. Suppose an unbiased preliminary bit $B$ becomes visible to the adversary before its final continue/abort decision. If it continues, output $B$; if it aborts, use a fresh independent unbiased fallback bit $C$. This grants the fallback an ideal property and isolates the abort decision.
\begin{proposition}
Continuing when $B=1$ and aborting when $B=0$ makes the probability of final output 1 equal to $3/4$. Three independent repetitions of this attacked experiment followed by majority have output-1 probability $27/32$.
\end{proposition}
\begin{proof}
The first probability is $\Pr(B=1)+\Pr(B=0,C=1)=1/2+1/4=3/4$, giving bias $1/4$. For independent repetitions the number of ones is binomial with parameters 3 and $3/4$. Its probability of being at least two is
\[
3(3/4)^2(1/4)+(3/4)^3=27/32,
\]
whose bias from $1/2$ is $11/32$, larger than $1/4$.
\end{proof}
Thus a fair fallback does not undo selective abort after seeing a correlated preliminary result, and naive repetition followed by majority can amplify a common directional bias. This is not an analysis of the source's carefully scheduled majority-style protocol, nor a lower bound for arbitrary $m$-round protocols. The adversary's information timing and the independence of repetitions are assumptions of this auxiliary experiment.
\subsection{Verification and remaining gaps}
All four pairs $(B,C)$ and all 64 triples of such pairs were enumerated, reproducing $3/4$ and $27/32$ exactly. The source's actual challenge of removing logarithmic factors under its stated assumptions remains untouched: no protocol with the requested bias and no general impossibility theorem is proved. The example only rejects a specific abort-handling and amplification route.

\section{References and source attribution}
Primary workshop report: \url{https://doi.org/10.4171/owr/2014/35}. Relevant passages read: 2014/35 pp. 1953--1954, round complexity, broadcast/abort model, three-party bias and the logarithmic-factor question.

All retrievals were checked on 11 October 2026. Original question attribution remains in the preserved Sources section. External theorems are attributed and not claimed as new discoveries or fully verified merely by citation.


\paragraph{Integration review.} Codex (GPT-6; exact serving revision not exposed) reviewed the written arguments and expanded layout and source-reference wrappers. Repository metadata and the production web renderer were checked using Python 3.13 and JavaScript with MathJax 3.2.2. This model-assisted review is not independent human proof verification; no PDF compilation or formal proof checking is asserted.

\section{Completion Estimate}
\noindent\textbf{COMPLETION: 0\%}
No substantiated new progress toward the entire original target is claimed. The worked results and failed route are delimited above; a known result, illustrative model, or auxiliary obstruction is not counted as original-target completion.

\section{Final status and attribution}
\textbf{Attempt status: unresolved.} This describes this research attempt, not a new ruling on the source catalogue's status or on external literature claims.
The mathematics and this self-audit were generated by OpenAI \texttt{gpt-6-astra}, reasoning \texttt{ultra}, on 11 October 2026. Codex (GPT-6; exact serving revision not exposed) orchestrated the batch, reviewed the written mathematical arguments, and contributed formatting/repository review. No human mathematical authorship, independent human proof checking, or formal verification is claimed. The historical preparation paragraphs in the reproduced catalogue body concern that body only.

Contribution prompt A was applied to the complete definition and supplied source record, followed by the adversarial checks of prompt B. Local reproduction files are \texttt{batch019/checks.py}, \texttt{computations.json}, and the per-ID audit JSON. They were not public links at preparation time. Numerical checks are finite and do not prove unrestricted results. No LaTeX engine was found locally; compilation was not performed by this research agent.

\end{document}
